\subsection{乘积测度的连续性}

命题 5. --- 对每个测度 \(\lambda_0 \in \mathcal{M}(\mathrm{X};\mathbf{C})\) ，映射 \(\mu \mapsto \lambda_0 \otimes \mu\) （ \(\mathcal{M}(\mathrm{Y};\mathbf{C})\) 到 \(\mathcal{M}(\mathrm{X} \times \mathrm{Y};\mathbf{C})\) 中的）是淡连续的。

对每个函数 \(f \in \mathcal{K}(\mathrm{X} \times \mathrm{Y}; \mathbf{C})\) ，我们知道函数 \(h(y) = \int f(x, y) d\lambda_0(x)\) 属于 \(\mathcal{K}(\mathrm{Y}; \mathbf{C})\) （No.~1, 引理 2），且 \(\langle f, \lambda_0 \otimes \mu \rangle = \langle h, \mu \rangle\) ，由此得该命题。

命题 6. --- 当 \(\mathcal{M}(\mathrm{X};\mathbf{C})\) 、 \(\mathcal{M}(\mathrm{Y};\mathbf{C})\) 与 \(\mathcal{M}(\mathrm{X}\times\mathrm{Y};\mathbf{C})\) 都赋以严格紧收敛拓扑（§1, No.~10）时，双线性映射 \((\lambda,\mu)\mapsto\lambda\otimes\mu\) （ \(\mathcal{M}(\mathrm{X};\mathbf{C})\times\mathcal{M}(\mathrm{Y};\mathbf{C})\) 到 \(\mathcal{M}(\mathrm{X}\times\mathrm{Y};\mathbf{C})\) 中的）对 \(\mathcal{M}(\mathrm{X};\mathbf{C})\) 与 \(\mathcal{M}(\mathrm{Y};\mathbf{C})\) 的淡有界子集全体是亚连续的（TVS, III, §5, No.~3）。

设 \(K \subset X\) ， \(L \subset Y\) 是两个紧集，A 是 \(\mathcal{K}(X \times Y, K \times L; C)\) 的一个紧子集，B 是 \(\mathcal{M}(X; C)\) 的一个淡有界闭子集；已知 B 是淡紧的（§1, No.~9, 命题 15），因而对严格紧收敛拓扑也是紧的（§1, No.~10, 命题 17）。另一方面，Banach 空间 \(\mathcal{K}(X \times Y, K \times L; C)\) 等距同构于 \(\mathcal{K}(X, K; \mathcal{K}(Y, L; C))\) （No.~1, 引理 1）；映射 \(\varphi\) （ \(\mathcal{K}(X, K; \mathcal{K}(Y, L; C)) \times \mathcal{M}(X; C)\) 到 \(\mathcal{K}(Y, L; C)\) 中的）——使 \(\varphi(g, \lambda)\) 为由 \(h(y) = \int g(x, y) d\lambda(x)\) 定义的函数 h——依据 §3, No.~4 的命题 8 与 9 是分别连续的。由于 \(\mathcal{K}(X, K; \mathcal{K}(Y, L; C))\) 是桶式的，由此可知映射 \(\varphi\) 相对于 \(\mathcal{M}(X; C)\) 的淡有界子集是亚连续的（TVS, III, §5, No.~3, 命题 6）；因此这一映射在 \(A \times B\) 上的限制是连续的（同处，命题 4）。于是 \(A \times B\) 在这一映射下的像 C 是 Banach 空间 \(\mathcal{K}(Y, L; C)\) 的一个紧子集。而 C 不是别的，正是函数 \(h(y) = \int f(x, y) d\lambda(x)\) （当 f 跑遍 A 且 \(\lambda\) 跑遍 B 时）所成之集；于是依据 No.~1 的公式 (3)，条件 \(\lambda \in B\) 与 \(\mu \in C^{\circ}\) 蕴涵 \(\lambda \otimes \mu \in A^{\circ}\) 。鉴于严格紧收敛拓扑的定义，这就证明了该命题（TVS, III, §5, No.~3, 定义 2）。

\pandocbounded{\includegraphics[keepaspectratio]{images/part001_f639fdbe.jpg}}

当把严格紧收敛拓扑换成淡拓扑时，命题 6 的结论不再成立（习题 2 c)）。然而，若 B（相应地 C）是 \(\mathcal{M}(\mathrm{X};\mathbf{C})\) （相应地 \(\mathcal{M}(\mathrm{Y};\mathbf{C})\) ）的一个淡有界子集，则 \(\mathrm{B}\times \mathrm{C}\) 在映射 \((\lambda ,\mu)\mapsto \lambda \otimes \mu\) 下的像在 \(\mathcal{M}(\mathrm{X}\times \mathrm{Y};\mathbf{C})\) 中是淡有界的，因而这一映射在 \(\mathrm{B}\times \mathrm{C}\) 上的限制是淡连续的，这依据命题 6、 §1, No.~10 的命题 17 以及 TVS, III, §5, No.~3 的命题 4（参见习题 3）。

